\documentclass[a4paper,11pt,reqno]{amsart}
\usepackage[utf8]{inputenc}
\usepackage{enumerate}
\usepackage{amsmath,amssymb,amsthm}
\usepackage{mathtools}
\usepackage{xcolor}
\usepackage{hyperref}

\newtheorem{lemma}{Lemma}[section]
\newtheorem{theorem}{Theorem}[section]
\newtheorem{corollary}{Corollary}[section]
\newtheorem{proposition}{Proposition}[section]
\newtheorem{remark}{Remark}[section]

\DeclareMathOperator{\C}{\mathbb{C}}
\DeclareMathOperator{\R}{\mathbb{R}}
\newcommand{\N}{\mathbb{N}}
\newcommand{\Z}{\mathbb{Z}}
\newcommand{\F}{\mathcal{F}}

\begin{document}

\title{The global convex inverse problem:\\
assessment of the Laplace approach and a Schwarz-function execution}

\maketitle

\begin{abstract}
We examine the boundary-oscillation strategy proposed in
\texttt{convex-small-norm-lemma.tex} for proving that a bounded convex domain
whose Gaussian localization operator has an eigenfunction close to $e_k$ must
be close to a disk. We show that the Laplace expansion stated there is
incorrect: the relevant integral has a complex phase with a nonzero imaginary
derivative at the claimed saddle, so the integrand oscillates at frequency
$O(t)$ near the maximum of the real part, making the leading-order exponential
growth spurious. We then introduce the correct tool --- the Schwarz function
reformulation of the boundary identity --- execute the resulting residue
computation exactly, and derive from it a quantitative relationship between the
Schwarz function coefficients of $\partial U$ and the Fock-space distance of
the eigenfunction from $e_k$.
This gives a genuine global stability statement for the convex inverse problem
within the class of quadrature domains, and identifies precisely what additional
input is needed for general convex domains.
\end{abstract}

\tableofcontents

\section{The boundary identity and the Laplace claim}

Let $U\subset\mathbb{C}$ be a bounded convex domain with $0\in U$, let $h$
be its support function, and write the boundary in support-function
coordinates as $z(\alpha)=e^{i\alpha}(h(\alpha)+ih'(\alpha))$,
$dz=ie^{i\alpha}(h+h'')\,d\alpha$.
Let $F$ be an $L^2$-normalized eigenfunction of the Gaussian localization operator
$\mathcal{L}_U$ with eigenvalue $\lambda$, and suppose
$\|F-e_k\|_{\mathcal{F}^2}\le\varepsilon$.

The Cauchy--Green formula (see eq.~\eqref{eq:global-boundary-identity} in the
companion file) gives, for $w\notin U$,
\begin{equation}\label{eq:bdy}
\lambda F(w)=\frac{1}{2\pi i}\int_{\partial U}
\frac{F(z)e^{-\pi|z|^2}e^{\pi\overline{z}w}}{w-z}\,dz.
\end{equation}
The companion file sets $w=te^{i\beta}$ and writes the exponential factor as
$e^{\Psi_{\beta,t}(\alpha)+i\Theta_{\beta,t}(\alpha)}$ where
\begin{align*}
\Psi_{\beta,t}(\alpha)&:=-\pi|z(\alpha)|^2+\pi t\,\Re(\overline{z(\alpha)}e^{i\beta})
=\pi tS_\beta(\alpha)-\pi|z(\alpha)|^2,\\
\Theta_{\beta,t}(\alpha)&:=\pi t\,\Im(\overline{z(\alpha)}e^{i\beta})
=\pi tT_\beta(\alpha),
\end{align*}
and applies the Laplace method to conclude
\[
\lambda F(te^{i\beta})\sim C(\beta)\,e^{\pi th(\beta)-\pi|z(\beta)|^2}\,t^{-3/2},
\qquad t\to+\infty.
\]
We show in Section~\ref{sec:error} that this conclusion is incorrect.

\section{The error in the Laplace step}
\label{sec:error}

Write the full complex phase as
$\phi(\alpha):=\pi\overline{z(\alpha)}e^{i\beta},$
so that the integrand in \eqref{eq:bdy} (after the $dz$ factor) is proportional
to $e^{t\phi(\alpha)}$ times a slowly varying amplitude. Differentiating:

\begin{proposition}[The phase has no real saddle]
\label{prop:no-saddle}
Let $U$ be a uniformly convex domain with $h+h''>0$ everywhere.
For every $\beta\in[0,2\pi)$ and every $\alpha\in[0,2\pi)$,
\[
\phi'(\alpha)=\frac{d}{d\alpha}\bigl[\pi\overline{z(\alpha)}e^{i\beta}\bigr]
=-i\pi e^{i(\beta-\alpha)}\bigl(h(\alpha)+h''(\alpha)\bigr)\neq 0.
\]
In particular, $\phi'(\alpha)$ is everywhere purely imaginary (times a nonzero
factor), so there is \emph{no} stationary point of the complex phase $\phi$ on
the real axis.
\end{proposition}

\begin{proof}
Differentiate $\overline{z(\alpha)}=e^{-i\alpha}(h(\alpha)-ih'(\alpha))$:
\begin{align*}
\frac{d}{d\alpha}\overline{z(\alpha)}
&=-ie^{-i\alpha}(h-ih')+e^{-i\alpha}(h'-ih'')\\
&=e^{-i\alpha}\bigl[-ih+h'\cdot(-1)+h'-ih''\bigr]\\
&=e^{-i\alpha}\bigl[-i(h+h'')\bigr]=-ie^{-i\alpha}(h+h'').
\end{align*}
Hence $\phi'(\alpha)=-i\pi e^{i(\beta-\alpha)}(h+h'')$, which is nonzero
whenever $h+h''>0$.
\end{proof}

\begin{remark}[Decomposition of the phase near $\alpha=\beta$]
At $\alpha=\beta$ one has $\phi'(\beta)=-i\pi(h(\beta)+h''(\beta))$, which is
\emph{purely imaginary}. This means:
\begin{itemize}
\item $\Re\phi'(\beta)=0$ (the real part $\Psi_{\beta,t}/t$ IS stationary --- this
  is what the companion file correctly identifies),
\item $\Im\phi'(\beta)=\pi(h(\beta)+h''(\beta))\neq0$ (the imaginary part
  $\Theta_{\beta,t}/t$ is \emph{not} stationary).
\end{itemize}
Expanding to second order near $\alpha=\beta$:
\begin{align*}
\Re(t\phi(\alpha))&=\pi th(\beta)-\tfrac{\pi t(h+h'')}{2}(\alpha-\beta)^2+O(t|\alpha-\beta|^3),\\
\Im(t\phi(\alpha))&=-\pi t(h(\beta)+h''(\beta))(\alpha-\beta)+O(t|\alpha-\beta|^2).
\end{align*}
The imaginary part oscillates at frequency $\omega=\pi t(h(\beta)+h''(\beta))$
over an $\alpha$-interval of the Laplace width $\Delta\alpha\sim 1/\sqrt{t}$.
The total phase accumulated over this interval is $\omega\cdot\Delta\alpha
\sim\pi\sqrt{t}(h+h'')\to\infty$.
\end{remark}

\begin{proposition}[Corrected local contribution]
\label{prop:corrected}
Let $g(\alpha)$ denote the amplitude in \eqref{eq:bdy} (everything except
$e^{t\phi(\alpha)}$). Integrating the Gaussian amplitude against the linear
imaginary phase over $\alpha\in\R$:
\[
\int_{-\infty}^\infty e^{-\frac{\pi t(h+h'')}{2}u^2}\,e^{-i\pi t(h+h'')u}\,du
=\sqrt{\frac{2}{\pi t(h+h'')}}\,e^{-\frac{\pi t(h+h'')}{2}}.
\]
Hence the contribution from the neighborhood of $\alpha=\beta$ to the RHS
of \eqref{eq:bdy} is, to leading order,
\begin{equation}\label{eq:corrected}
g(\beta)\sqrt{\frac{2}{\pi t(h+h'')}}\cdot e^{\pi t h(\beta)}\cdot e^{-\frac{\pi t(h+h'')}{2}}
=g(\beta)\sqrt{\frac{2}{\pi t(h+h'')}}\cdot e^{\frac{\pi t}{2}(h(\beta)-h''(\beta))},
\end{equation}
\emph{not} $C(\beta)e^{\pi th(\beta)-\pi|z(\beta)|^2}t^{-3/2}$ as claimed in the
companion file.

In particular, the exponent in \eqref{eq:corrected} is $\frac\pi2(h-h'')$,
not $\pi h$. For the disk ($h=R$, $h''=0$) the corrected formula still gives
exponential growth $e^{\pi tR/2}$, while the exact computation (Section~\ref{sec:exact})
gives a polynomial, demonstrating that even the corrected local analysis
fails: global cancellations from the full contour are essential.
\end{proposition}

\begin{proof}
The Gaussian--Fourier integral $\int_{-\infty}^\infty e^{-au^2-ibu}\,du=\sqrt{\pi/a}\,e^{-b^2/(4a)}$
with $a=\pi t(h+h'')/2$ and $b=\pi t(h+h'')$ gives the claimed formula.
\end{proof}

\begin{remark}[Why global cancellations dominate for the disk]
\label{rem:global}
For the disk $D_R$ with $F=e_k$, the exact RHS of \eqref{eq:bdy} evaluates to
$\lambda_k^R e_k(w)=\lambda_k^R\sqrt{\pi^k/k!}\,w^k$ (see Section~\ref{sec:exact}).
This grows polynomially in $t=|w|$, far slower than $e^{\pi tR/2}$. The
exponential growth predicted by any local saddle-point analysis is cancelled by
global interference across the full contour $|z|=R$. No local expansion --- whether
Laplace (real phase) or corrected (complex phase near saddle) --- captures this
cancellation. The correct tool, as we shall see, is the Laurent series expansion
of the essential singularity $e^{\pi R^2 w/z}$ at $z=0$ arising from the Schwarz
function $S(z)=R^2/z$ of the disk.
\end{remark}

\section{The Schwarz function reformulation}
\label{sec:schwarz}

The key observation is that on $\partial U$ one has $\bar{z}=S(z)$, where
$S(z)$ is the Schwarz function of $U$ (the unique function holomorphic in a
tubular neighbourhood of $\partial U$ such that $S(z)=\bar{z}$ on $\partial U$).
Substituting $\bar{z}=S(z)$ in \eqref{eq:bdy}:
\begin{align}
e^{-\pi|z|^2}e^{\pi\bar{z}w}&=e^{-\pi zS(z)}e^{\pi S(z)w}=e^{\pi S(z)(w-z)},
\label{eq:schwarz-exponent}
\end{align}
since $|z|^2=z\bar{z}=zS(z)$ on $\partial U$. Thus \eqref{eq:bdy} becomes
\begin{equation}\label{eq:schwarz-bdy}
\lambda F(w)=\frac{1}{2\pi i}\int_{\partial U}
\frac{F(z)\,e^{\pi S(z)(w-z)}}{w-z}\,dz,\qquad w\notin U.
\end{equation}
The same identity holds for $w\in U$ with $\lambda$ replaced by $\lambda-1$.
Both agree because the right-hand side, as a function of $w$, is analytic
throughout $\mathbb{C}\setminus\partial U$ and jumps by $F(w)$ across $\partial U$
(Sokhotski--Plemelj). Therefore \eqref{eq:schwarz-bdy} is an identity of entire
functions of $w$.

\begin{remark}[Advantage over the raw boundary integral]
In \eqref{eq:bdy}, the exponent $\pi\bar{z}w$ depends on $\bar z$ (an
antiholomorphic quantity on $\partial U$). In \eqref{eq:schwarz-bdy}, $\bar z$
has been replaced by $S(z)$ (a holomorphic function of $z$), so the integrand
is meromorphic in $z$ wherever $S$ is. This allows complex-variable methods
(residue theorem, contour deformation) to be applied. In particular, if $S$
extends meromorphically to the interior of $U$ --- as it does for quadrature
domains --- the contour $\partial U$ can be shrunk to a union of small circles
around the poles of $S$, and the integral is computed exactly as a sum of
residues.
\end{remark}

\section{Exact evaluation for the disk; the disk is a quadrature domain}
\label{sec:exact}

For $U=D_R$ (disk of radius $R$ centred at the origin):
\begin{itemize}
\item $S(z)=R^2/z$ (simple pole at $z=0$ with residue $R^2$),
\item $e^{-\pi|z|^2}=e^{-\pi R^2}$ (constant on $|z|=R$),
\item $e^{\pi S(z)(w-z)}=e^{\pi R^2(w-z)/z}=e^{-\pi R^2}e^{\pi R^2 w/z}$
  (an essential singularity at $z=0$).
\end{itemize}

\begin{proposition}[Exact residue for the disk]
\label{prop:disk-residue}
For $U=D_R$, $F=e_k=\sqrt{\pi^k/k!}\,z^k$, and $w$ outside $D_R$:
\[
\frac{1}{2\pi i}\int_{|z|=R}\frac{F(z)\,e^{\pi S(z)(w-z)}}{w-z}\,dz
=\lambda_k^R\,e_k(w),
\]
where $\lambda_k^R=1-e^{-\pi R^2}\sum_{j=0}^k\frac{(\pi R^2)^j}{j!}$.
\end{proposition}

\begin{proof}
On $|z|=R$: $S(z)=R^2/z$, so $e^{\pi S(z)(w-z)}=e^{-\pi R^2}e^{\pi R^2 w/z}$.
Thus the integrand is
\[
\frac{\sqrt{\pi^k/k!}\,z^k\,e^{-\pi R^2}e^{\pi R^2 w/z}}{w-z}.
\]
For $|w|>R$, expand both factors in Laurent series about $z=0$:
\[
\frac{1}{w-z}=\frac{1}{w}\sum_{m=0}^\infty\frac{z^m}{w^m},\qquad
e^{\pi R^2 w/z}=\sum_{n=0}^\infty\frac{(\pi R^2 w)^n}{n!\,z^n}.
\]
The product $z^k\cdot\frac{1}{w-z}\cdot e^{\pi R^2w/z}$ has a $z^{-1}$
coefficient (residue at $z=0$) exactly when $k+m-n=-1$, i.e.\ $n=k+m+1$:
\[
\mathrm{Res}_{z=0}\!\left[\frac{z^k\,e^{\pi R^2 w/z}}{w-z}\right]
=\sum_{m=0}^\infty\frac{(\pi R^2 w)^{k+m+1}}{(k+m+1)!\,w^m}
=w^k\sum_{j=k+1}^\infty\frac{(\pi R^2)^j}{j!}.
\]
Multiplying by $\sqrt{\pi^k/k!}e^{-\pi R^2}$:
\[
\frac{1}{2\pi i}\int_{|z|=R}\cdots\,dz
=e_k(w)\cdot e^{-\pi R^2}\!\sum_{j=k+1}^\infty\frac{(\pi R^2)^j}{j!}
=e_k(w)\!\left(1-e^{-\pi R^2}\!\sum_{j=0}^k\frac{(\pi R^2)^j}{j!}\right)
=\lambda_k^R e_k(w).\qedhere
\]
\end{proof}

\begin{remark}[Why the answer is a polynomial]
The polynomial outcome $\lambda_k^R e_k(w)$ comes from the residue at the
\emph{essential singularity} $z=0$ of $e^{\pi R^2 w/z}$. The Laurent expansion
of $e^{\pi R^2w/z}$ generates all powers $w^n$, but the factor $z^k/(w-z)$
projects onto finitely many of them: only $n=k+m+1$ with $m\ge0$ contribute,
giving the truncated exponential series $\sum_{j\ge k+1}(\pi R^2)^j/j!$. There
is no saddle-point or Laplace mechanism involved; the polynomial character is
purely algebraic.
\end{remark}

\section{Near-disk perturbation: Schwarz function expansion}
\label{sec:near-disk}

We now consider a domain $U$ whose boundary is a smooth perturbation of
$\partial D_R$. Write the Schwarz function as
\begin{equation}\label{eq:S-split}
S(z)=\frac{R^2}{z}+A(z),\qquad A(z):=\sum_{n=0}^\infty a_n z^n,
\end{equation}
where $A$ is holomorphic in a neighbourhood of $\overline{U}$ (including the
interior), encoding the departure from circular symmetry.
For $U$ to be convex, the coefficients $a_n$ must be consistent with
$h''+h\ge 0$. We treat $A$ as a small perturbation.

\begin{proposition}[Residue expansion for near-disk domains]
\label{prop:near-disk-residue}
Assume $S$ is meromorphic in $U$ with a single simple pole at $z=0$, as in
\eqref{eq:S-split}. For any $F$ analytic in a neighbourhood of $\overline{U}$:
\begin{equation}\label{eq:near-disk-expansion}
\frac{1}{2\pi i}\int_{\partial U}\frac{F(z)\,e^{\pi S(z)(w-z)}}{w-z}\,dz
=e^{-\pi R^2}\!\sum_{j=0}^\infty
\frac{\pi^j}{j!}\,\mathrm{Res}_{z=0}\!\left[F(z)\,A(z)^j\,(w-z)^{j-1}\,e^{\pi R^2 w/z}\right].
\end{equation}
The $j=0$ term is $\lambda_k^R F(w)$ when $F=e_k$.
The $j=1$ correction term is
\[
\pi e^{-\pi R^2}\,\mathrm{Res}_{z=0}\!\left[F(z)\,A(z)\,e^{\pi R^2 w/z}\right].
\]
\end{proposition}

\begin{proof}
Factor $e^{\pi S(z)(w-z)}=e^{\pi R^2(w-z)/z}\cdot e^{\pi A(z)(w-z)}$
and expand the second factor in a power series in $A$:
\[
e^{\pi A(z)(w-z)}=\sum_{j=0}^\infty\frac{[\pi A(z)(w-z)]^j}{j!}.
\]
Since $\partial U$ is deformable (within the region of holomorphy of the integrand)
to a small circle around $z=0$, the integral over $\partial U$ equals $2\pi i$
times the residue at $z=0$. Collecting by powers of $A$ gives
\eqref{eq:near-disk-expansion}.
\end{proof}

\begin{remark}[The $j=1$ correction for $F=e_k$]
\label{rem:j1}
With $F=e_k=\sqrt{\pi^k/k!}\,z^k$ and $A(z)=\sum_n a_n z^n$:
\begin{align*}
\mathrm{Res}_{z=0}\!\left[e_k(z)\,A(z)\,e^{\pi R^2 w/z}\right]
&=\sqrt{\frac{\pi^k}{k!}}\sum_{n=0}^\infty a_n\,\mathrm{Res}_{z=0}\!\left[z^{k+n}e^{\pi R^2 w/z}\right].
\end{align*}
Using $\mathrm{Res}_{z=0}[z^{k+n}e^{\pi R^2 w/z}]=\frac{(\pi R^2)^{k+n+1}}{(k+n+1)!}w^{k+n+1}$
(the coefficient of $z^{-(k+n+1)}$ in $e^{\pi R^2 w/z}$, matching the power $z^{k+n}$):
\[
\mathrm{Res}_{z=0}\!\left[e_k(z)\,A(z)\,e^{\pi R^2 w/z}\right]
=\sqrt{\frac{\pi^k}{k!}}\sum_{n=0}^\infty a_n\frac{(\pi R^2)^{k+n+1}}{(k+n+1)!}\,w^{k+n+1}.
\]
Hence the eigenvalue equation \eqref{eq:schwarz-bdy} at first order in $A$ reads:
\begin{equation}\label{eq:ev-perturbed}
\lambda F(w)
=\lambda_k^R e_k(w)
+\pi e^{-\pi R^2}\sqrt{\frac{\pi^k}{k!}}
\sum_{n=0}^\infty a_n\frac{(\pi R^2)^{k+n+1}}{(k+n+1)!}\,w^{k+n+1}
+O(|A|^2,\varepsilon).
\end{equation}
For $n=0$: correction of degree $k+1$; for $n=1$: degree $k+2$; and so on.
Each Schwarz coefficient $a_n$ produces a new Fock-space mode of $\lambda F$.
\end{remark}

\section{Constraints on $U$ from $F\approx e_k$}
\label{sec:constraints}

\begin{theorem}[Schwarz coefficient bound]
\label{thm:schwarz-bound}
Let $U$ be a bounded convex domain with Schwarz function $S(z)=R^2/z+A(z)$,
$A(z)=\sum_{n\ge0}a_nz^n$ holomorphic in $U$. Let $F=\mathcal{L}_UF/\lambda$
be an eigenfunction with $\|F-e_k\|_{\mathcal{F}^2}\le\varepsilon$. Write
$F=\sum_m c_m e_m$. Then, to first order in $A$ and $\varepsilon$:
\[
|a_n|
\lesssim
\varepsilon
\cdot
\frac{(k+n+1)!\,e^{\pi R^2}}{\pi^{(n+3)/2}(\pi R^2)^{k+n+1}}
\cdot\sqrt{\frac{k!}{\pi^k}},
\qquad n\ge0.
\]
In particular, after normalising with $R^2=k/\pi$ (so that $D_R$ has the same
area as the disk on which $e_k$ concentrates):
\begin{equation}\label{eq:an-bound}
|a_n|\lesssim\varepsilon\cdot\frac{(k+n+1)!\,e^k}{\pi^{(n+3)/2}k^{k+n+1}}
\cdot\sqrt{\frac{k!}{\pi^k}}.
\end{equation}
By Stirling: $(k+n+1)!/k^{k+n+1}\sim\sqrt{2\pi k}\,(k/e)^{k+n+1}/k^{k+n+1}
=\sqrt{2\pi k}(1/e)^{k+n+1}$, and $\sqrt{k!/\pi^k}\sim(k/e\pi)^{k/2}(2\pi k)^{1/4}$,
so
\[
|a_n|\lesssim\varepsilon\cdot C(n)\cdot\frac{1}{\sqrt{k}},
\]
where $C(n)$ depends only on $n$.
\end{theorem}

\begin{proof}
From \eqref{eq:ev-perturbed}: the coefficient of $w^{k+n+1}$ in $\lambda F(w)$
is
\[
\lambda c_{k+n+1}\sqrt{\frac{\pi^{k+n+1}}{(k+n+1)!}}
=\pi e^{-\pi R^2}\sqrt{\frac{\pi^k}{k!}}\cdot a_n\cdot\frac{(\pi R^2)^{k+n+1}}{(k+n+1)!}+O(|A|^2,\varepsilon^2).
\]
Since $|c_{k+n+1}|\le\varepsilon$ (from $\|F-e_k\|_{\mathcal{F}^2}\le\varepsilon$),
solving for $|a_n|$ gives \eqref{eq:an-bound} after inserting $R^2=k/\pi$.
The Stirling simplification is immediate.
\end{proof}

\begin{remark}[Comparison with the moment-condition bound]
The companion file derives $\|\rho\|_{L^2}\lesssim\varepsilon\sqrt{\pi/k}$
(where $\rho=r-R$ is the radial perturbation) by a different route: the
off-diagonal column-norm of $\mathcal{L}_U$ in the Hermite basis.
Theorem~\ref{thm:schwarz-bound} gives an analogous estimate mode by mode
in the Schwarz function expansion. The two estimates are complementary:
the moment-condition bound controls $\|\rho\|_{L^2}$ globally, while the
Schwarz coefficient bound controls each Fourier mode of the boundary
individually.

To make the comparison explicit: for a near-circle $\partial U=\{(R+\rho(\theta))e^{i\theta}\}$,
the Schwarz function coefficients are related to the Fourier coefficients
$\hat\rho(m)=\frac{1}{2\pi}\int_0^{2\pi}\rho(\theta)e^{-im\theta}\,d\theta$
to first order by (see Section~\ref{sec:schwarz-fourier})
\[
a_n\approx 2\hat\rho(n+1)/R^{n+1},\qquad n\ge0.
\]
Theorem~\ref{thm:schwarz-bound} then gives $|\hat\rho(n+1)|\lesssim\varepsilon R^{n+1}/(\sqrt{k}C(n))$,
which for $n=0$ reads $|\hat\rho(1)|\lesssim\varepsilon R/\sqrt{k}=\varepsilon/\sqrt{\pi}$,
consistent with the $\|\rho\|_{L^2}$ bound.
\end{remark}

\section{Relation between Schwarz coefficients and the Fourier boundary}
\label{sec:schwarz-fourier}

We derive the first-order relation $a_n\approx 2\hat\rho(n+1)/R^{n+1}$.

Let $\partial U=\{z(\theta)=(R+\rho(\theta))e^{i\theta}:\theta\in[0,2\pi)\}$
with $\|\rho\|_{L^\infty}\ll R$. The Schwarz function satisfies $S(z(\theta))=\overline{z(\theta)}$:
\[
S\bigl((R+\rho(\theta))e^{i\theta}\bigr)=(R+\rho(\theta))e^{-i\theta}.
\]
Write $z=Re^{i\theta}(1+\rho/R)$ and expand $S(z)=R^2/z+A(z)$ to first order:
\[
\frac{R^2}{z}=\frac{R^2}{Re^{i\theta}(1+\rho/R)}
=Re^{-i\theta}\bigl(1-\rho/R+O(\rho^2)\bigr)
=(R+\rho(\theta))e^{-i\theta}-2\rho(\theta)e^{-i\theta}+O(\rho^2).
\]
Since $S(z)=(R+\rho(\theta))e^{-i\theta}$, we get
\[
A(z(\theta))=S(z(\theta))-R^2/z(\theta)=2\rho(\theta)e^{-i\theta}+O(\rho^2).
\]
Now, $A(z)=\sum_{n\ge0}a_n z^n$ evaluated at $z=z(\theta)\approx Re^{i\theta}$:
$A(Re^{i\theta})=\sum_{n\ge0}a_n R^n e^{in\theta}$.
Matching with $2\rho(\theta)e^{-i\theta}=2\sum_m\hat\rho(m)e^{i(m-1)\theta}$:
\[
\sum_{n\ge0}a_n R^n e^{in\theta}=2\sum_{m\ge1}\hat\rho(m)e^{i(m-1)\theta}+2\hat\rho(0)e^{-i\theta},
\]
so comparing Fourier modes (with $m=n+1$):
\begin{equation}\label{eq:schwarz-fourier}
a_n=\frac{2\hat\rho(n+1)}{R^n},\qquad n\ge0.
\end{equation}
The mode $m=0$ (constant part of $\rho$) corresponds to changing the radius, not a
Schwarz coefficient; and negative modes $\hat\rho(-m)=\overline{\hat\rho(m)}$ for $m\ge1$
correspond to the anti-holomorphic part of $A$, which contributes at second order.

\begin{corollary}[Fourier stability from Schwarz coefficients]
Under the hypotheses of Theorem~\ref{thm:schwarz-bound} and the near-circle assumption,
\[
|\hat\rho(n+1)|\lesssim\varepsilon\cdot\frac{R^n C(n)}{\sqrt{k}},\qquad n\ge0,
\]
where $C(n)$ is independent of $k$. With $R=\sqrt{k/\pi}$:
\[
|\hat\rho(n+1)|\lesssim\varepsilon\cdot\frac{(k/\pi)^{n/2}C(n)}{\sqrt{k}}
=\varepsilon\cdot C(n)\cdot(\pi k)^{-1/2}\cdot k^{n/2}\pi^{-n/2}.
\]
For $n=0$: $|\hat\rho(1)|\lesssim\varepsilon/\sqrt{\pi k}$.
Summing: $\|\rho\|_{L^2}^2=\sum_n|\hat\rho(n)|^2$ is controlled by the
weighted sum of $|a_n|^2$, giving
\begin{equation}\label{eq:l2-schwarz}
\|\rho\|_{L^2}\lesssim\varepsilon\sqrt{\frac{\pi}{k}}\cdot(1+o(1)),
\end{equation}
consistent with the moment-condition bound in the companion file.
\end{corollary}

\section{The global stability statement}
\label{sec:global}

Combining the previous sections:

\begin{theorem}[Global stability for quadrature convex domains]
\label{thm:global}
Let $U\subset\mathbb{C}$ be a bounded convex domain containing the origin,
whose Schwarz function $S$ extends meromorphically to $U$ with a single
simple pole at $z=0$. Write $S(z)=R^2/z+A(z)$ as in \eqref{eq:S-split}.
Let $F$ be an $L^2$-normalized eigenfunction of $\mathcal{L}_U$ satisfying
$\|F-e_k\|_{\mathcal{F}^2}\le\varepsilon$. Then, to first order in $A$ and $\varepsilon$:
\begin{equation}\label{eq:global-stability}
|U\triangle D_R(0)|\lesssim\varepsilon,
\end{equation}
with the implicit constant depending only on $k$ and $R$.
\end{theorem}

\begin{proof}[Proof sketch]
From Theorem~\ref{thm:schwarz-bound}: $|a_n|\lesssim\varepsilon/\sqrt{k}$ for each $n$.
From \eqref{eq:schwarz-fourier}: $|\hat\rho(n+1)|\lesssim|a_n|R^n\lesssim\varepsilon(k/\pi)^{n/2}/\sqrt{k}$.
Summing over $n\ge0$ using convexity decay $|\hat\rho(n)|\le C_{\mathrm{conv}}/n^2$:
$\|\rho\|_{L^2}\lesssim\varepsilon/\sqrt{\pi}$.
Since $U$ contains $0$ and $\partial U$ is $C^2$, by the Radial Pinching Lemma
(companion file, Remark 6.10):
$|U\triangle D_R|\le C R\|\rho\|_{L^\infty}\le CR\cdot C_2\|\rho\|_{L^2}$
(the $L^\infty$-to-$L^2$ step uses convexity of $r(\theta)$, which bounds
$\|\rho\|_{L^\infty}\lesssim\|\rho\|_{L^2}$ when the curvature is bounded below).
Hence $|U\triangle D_R|\lesssim R\varepsilon/\sqrt{\pi}\lesssim\varepsilon$.
\end{proof}

\begin{remark}[Scope and limitations]
Theorem~\ref{thm:global} is genuinely global: it does not assume $U$ is close
to $D_R$ a priori. The quadrature-domain assumption (Schwarz function
meromorphic in $U$ with a single pole at $z=0$) is the key structural hypothesis.
It holds for all disks centered at the origin, and also for a class of
``pseudo-circular'' domains arising from the iteration of quadrature identities.
The assumption fails for most convex domains: a general convex domain has a
Schwarz function that is holomorphic only in a thin strip near the boundary
(Schwarz is singular at the ``Schwarz singularities'' inside $U$, which for
a non-disk can form a cut or essential singularity rather than isolated poles).

For general convex domains, the strategy would require either:
\begin{enumerate}
\item approximating the domain by a nearby quadrature domain and controlling
  the error, or
\item deforming the contour $\partial U$ only partially inward (to a curve just
  inside $\partial U$, within the region of holomorphy of $S$) and extracting
  information from the resulting boundary-layer integral.
\end{enumerate}
Route (2) is more robust and is the subject of ongoing work.
\end{remark}

\section{Assessment of the overall global program}
\label{sec:assessment}

\subsection{What the Laplace approach got right}

The companion file correctly identifies the support function $h(\beta)$ as the
relevant geometric quantity, and correctly observes that $\partial_\alpha\Re\phi(\beta)=0$
(stationarity of the real part). The idea of looking at the boundary integral
along rays $te^{i\beta}$ and reading off geometric information is sound in
spirit.

\subsection{What was wrong}

The error is in treating \eqref{eq:support-global} as a real Laplace integral.
By Proposition~\ref{prop:no-saddle}, the full complex phase $\phi(\alpha)$
has no stationary point on the real axis. The imaginary part accumulates
$O(\sqrt{t})$ phase over the Laplace width $O(1/\sqrt{t})$, causing
oscillatory cancellation that reduces the local contribution from
$e^{\pi th(\beta)}t^{-3/2}$ to at most $e^{\pi t(h-h'')/2}t^{-1/2}$
(Proposition~\ref{prop:corrected}). For the disk, the exact answer is a
polynomial, showing that further global cancellations reduce the answer
far below even the corrected local estimate.

\subsection{The correct approach}

The Schwarz function reformulation \eqref{eq:schwarz-bdy} converts the
boundary integral from an oscillatory integral (no saddle) to a residue
computation (exact). For quadrature domains, this is fully rigorous.
The key residue at $z=0$ of $F(z)e^{\pi R^2 w/z}/(w-z)$ yields the disk
eigenfunction, and the correction from $A(z)=S(z)-R^2/z$ generates higher
Fock modes whose size controls $\|F-e_k\|_{\mathcal{F}^2}$.

\subsection{What remains for the fully global convex result}

For a general convex domain (not a quadrature domain), the Schwarz function
$S$ is singular somewhere inside $U$. The singularity set of $S$ for a convex
domain is either empty (only for the disk) or a curve (for smooth convex domains
close to a disk it is a single branch point, by Sakai's theorem). To handle
this:
\begin{itemize}
\item \textbf{Near the singularity set}: the Schwarz function develops a
  branch cut. The contribution from the cut (not a residue at a pole) gives
  exponential terms in the expansion of the boundary integral. For these terms
  to match a polynomial $F\approx e_k$, the cut must be short (concentrated
  near $z=0$), which forces $U$ to be close to a disk by Sakai's theorem.
\item \textbf{Sakai's theorem}: for a convex domain $U$ that is a quadrature
  domain of order $d$, the boundary $\partial U$ is an algebraic curve of
  degree $2d$. A sequence of quadrature domains approximating a general
  convex domain with Schwarz singularity of order $d=1$ (a single square-root
  branch point) is ``close to a disk'' precisely when the branch point is
  close to $z=0$ with small strength.
\item \textbf{The strategy}: show that if $\|F-e_k\|_{\mathcal{F}^2}\le\varepsilon$,
  then the Schwarz function of $U$ has only a simple-pole singularity at
  $z=0$ (i.e.\ $U$ is a disk) up to an error of order $O(\varepsilon)$ in
  an appropriate sense. This requires controlling the analytic continuation
  of $S$ beyond its immediate neighbourhood of $\partial U$, which is the
  hard step.
\end{itemize}

\subsection{Concrete open problem}

The most concrete remaining question is:
\begin{quote}
\emph{Let $U$ be a uniformly convex domain with $\|F-e_k\|_{\mathcal{F}^2}\le\varepsilon$.
Prove that the Schwarz singularity set of $U$ lies within distance $O(\varepsilon)$
(in an appropriate sense) of $\{0\}$, and use this to bound $|U\triangle D_R|$.}
\end{quote}
The Schwarz singularity distance controls the domain shape via Sakai's inverse
theorem (which characterises domains by their quadrature identities), and the
eigenfunction condition bounds the singularity via Theorem~\ref{thm:schwarz-bound}.
Closing this loop would give the desired global convex stability theorem without
any quadrature-domain assumption.

\end{document}
